\documentclass[12pt]{article}

\usepackage{amssymb,amsmath,amsthm}

% gaya teorema (saya ingin semuanya memakai pencacah yang sama)
\theoremstyle{plain}
\newtheorem{theorem}{Teorema}[section]
\newtheorem{lemma}[theorem]{Lema}
\newtheorem{proposition}[theorem]{Proposisi}
\newtheorem{corollary}[theorem]{Korolari}
\newtheorem{conjecture}[theorem]{Konjektur}
\newtheorem{exercise}[theorem]{Latihan}
\theoremstyle{definition}
\newtheorem{definition}[theorem]{Definisi}
\newtheorem{example}[theorem]{Contoh}
\theoremstyle{remark}
\newtheorem{remark}[theorem]{Catatan}

% beberapa pengaturan penomoran
\numberwithin{equation}{section}

\usepackage{fancyhdr}
\usepackage{lastpage}
\setlength{\headheight}{15.2pt}
\renewcommand{\footrulewidth}{0.4pt}% nilai bakunya 0pt
\setlength{\footskip}{30pt}
\pagestyle{fancy}

\makeatletter
\let\ps@plain\ps@fancy 
\makeatother

\lhead{MATH 742}
\chead{Pengantar Teori Representasi}
\rhead{Semester Musim Semi 2023}
\lfoot{Terakhir direvisi: \today}
\cfoot{}
\rfoot{\thepage\ dari \pageref{LastPage} }

\begin{document}

%%%%%%%%%%%%%%%%%%
%%%%%%%%%%%%%%%%%%
\title{Pengantar Teori Representasi}
\author{Alexander Duncan}

\maketitle

\section{Konsep Dasar}

Beberapa konsep yang dibahas di sini juga dapat ditemukan dalam
\cite[\S{1}]{Serre}.
Kita akan menggunakan banyak sumber lain untuk materi ini,
tetapi kebanyakan sumber tersebut (dengan tepat!) memakai sudut pandang teori modul,
yang untuk sementara ingin kita tunda.

\begin{definition}
Misalkan $V$ suatu ruang vektor di atas medan $k$.
\emph{Grup linear umum dari $V$}, yang dinotasikan dengan $\operatorname{GL}(V)$, adalah grup
transformasi linear invertibel $V \to V$.
Untuk bilangan bulat nonnegatif $n$, kita memakai notasi
$\operatorname{GL}_n(k) := \operatorname{GL}(k^n)$.
\end{definition}

Grup $\operatorname{GL}_n(k)$ dapat diidentifikasi secara kanonik dengan himpunan matriks berukuran $n
\times n$ yang invertibel dan entri-entrinya berada dalam $k$.

\begin{definition}
Misalkan $G$ suatu grup. Suatu \emph{representasi linear dari $G$} adalah
homomorfisme grup $\rho : G \to \operatorname{GL}(V)$ untuk suatu ruang vektor $V$.
\emph{Derajat} suatu representasi linear adalah dimensi
ruang vektor $V$ yang bersesuaian.
Representasi linear $\rho$ disebut \emph{setia} jika $\rho$ injektif
sebagai pemetaan himpunan.
\end{definition}

Kita kerap menyebut “representasi $V$”, dengan menunjuk pada
ruang vektor yang mendasarinya, bukan pada homomorfisme $\rho$.
Kadang-kadang kita menulis $\rho_g$ alih-alih $\rho(g)$ karena notasi tersebut lebih mudah
dibaca. Kelak, kita cukup menulis “$g$” untuk $\rho(g)$ apabila
tidak ada risiko kebingungan.

\begin{example}
Untuk setiap grup $G$,
\emph{representasi nol} adalah representasi berderajat
$0$ (representasi ini tunggal hingga isomorfisme).
\end{example}

\begin{example}
Untuk setiap grup $G$,
\emph{representasi satuan} atau \emph{representasi trivial} adalah representasi
\[ \rho: G \to \operatorname{GL}_1(k) \]
yang diberikan oleh
\[
\rho(g)=\operatorname{id}_k
\]
untuk semua $g \in G$.
\end{example}

\begin{example}
Misalkan $G = \langle g \rangle$ grup siklik berorde $n$.
Terdapat representasi
\[ \rho : G \to \operatorname{GL}_1(\mathbb{C}) \]
yang didefinisikan oleh
\[
\rho(g^k)=e^{\frac{2\pi i}{n}k} \ .
\]
\end{example}

\begin{example}
Misalkan $D_8 = \langle s,r \mid s^2, r^4, (sr)^2 \rangle$
grup dihedral berorde $8$.
Terdapat representasi
\[ \rho : G \to \operatorname{GL}_2(k) \]
yang didefinisikan pada para pembangkit oleh
\[
\rho(s)=\begin{pmatrix} 0&1\\1&0 \end{pmatrix} ,\quad
\rho(r)=\begin{pmatrix} 0&-1\\1&0 \end{pmatrix}
\]
dan diperluas ke elemen-elemen lain dengan menggunakan fakta bahwa $\rho$
harus merupakan homomorfisme grup.
Perhatikan bahwa $\rho$ merupakan representasi \emph{setia} jika dan hanya jika
$\operatorname{char}(k) \ne 2$.
\end{example}


\begin{example}
Untuk setiap grup $G$ yang beraksi dari kiri pada suatu himpunan $X$,
misalkan $V$ ruang vektor dengan basis yang diindeks oleh $X$:
\[
V = \mathrm{span} \{ e_x : x \in X \} \ .
\]
Definisikan \emph{representasi permutasi} yang bersesuaian dengan $X$
sebagai representasi
\[
\rho : G \to \operatorname{GL}(V)
\]
yang didefinisikan sedemikian sehingga
\[
\rho(g)(e_x) = e_{g(x)}
\]
untuk semua $g \in G$ dan $x \in X$.
\end{example}

\begin{example}
Untuk setiap grup $G$, \emph{representasi reguler dari $G$}
adalah representasi permutasi dari aksi $G$ pada dirinya sendiri melalui
perkalian kiri.
Dengan kata lain, 
$V$ adalah ruang vektor dengan basis yang diindeks oleh $G$:
\[
V = \mathrm{span} \{ e_g : g \in G \}
\]
dan representasi
\[
\rho : G \to \operatorname{GL}(V)
\]
didefinisikan sedemikian sehingga
\[
\rho(g)(e_h) = e_{gh}
\]
untuk semua $g,h \in G$.
\end{example}

\begin{example} \label{ex:lin_alg_is_Z_rep_theory}
Representasi linear $\rho : \mathbb{Z} \to \operatorname{GL}_n(k)$
ditentukan sepenuhnya oleh citra $\rho(1)$
karena $\rho(n)=\rho(1)^n$ untuk semua $n \in \mathbb{Z}$.
Ini memberikan bijeksi kanonik antara representasi linear
$\mathbb{Z}$ berderajat $n$ dan
matriks invertibel berukuran $n \times n$.
\end{example}

\subsection{Pemetaan Ekuivarian}

\begin{definition}
Misalkan $V$ dan $W$ ruang vektor di atas medan yang sama.
Andaikan
\begin{align*}
\rho &: G \to \operatorname{GL}(V), \textrm{ dan}\\
\rho' &: G \to \operatorname{GL}(W)
\end{align*}
merupakan representasi linear dari grup $G$ yang sama.
Suatu pemetaan linear $\tau : V \to W$ disebut \emph{$G$-ekuivarian} jika
\[
\rho'(g) \circ \tau = \tau \circ \rho(g)
\]
untuk semua $g \in G$.
Kedua representasi tersebut disebut \emph{serupa} atau \emph{isomorfik} jika
terdapat pemetaan linear $G$-ekuivarian yang invertibel, $\tau : V \to W$.
Misalkan $\operatorname{Hom}_k^G(V,W)$ himpunan bagian dari
pemetaan linear $G$-ekuivarian $f : V \to W$.
\end{definition}

Istilah lama untuk pemetaan $G$-ekuivarian adalah “operator penjalin”,
tetapi istilah ini mulai menghilang (dan menurut saya, untunglah demikian).

\begin{example}
Jika $A$ dan $B$ merupakan matriks invertibel yang serupa, maka terdapat
matriks invertibel $P$ sedemikian sehingga $A=PBP^{-1}$.
Ingat bahwa $A^n=PB^nP^{-1}$ untuk semua bilangan bulat $n$.
Berdasarkan Contoh~\ref{ex:lin_alg_is_Z_rep_theory}, kita menyimpulkan bahwa
kelas ekuivalensi representasi linear
$\mathbb{Z}$ berderajat $n$ berkorespondensi secara bijektif kanonik dengan
kelas keserupaan matriks invertibel berukuran $n \times n$.
\end{example}

\begin{example}
Misalkan $\rho : \mathbb{Z}/4\mathbb{Z} \to
\operatorname{GL}_1(\mathbb{C})$
representasi dengan $\rho(1)=(i)$,
dan misalkan $\sigma : \mathbb{Z}/4\mathbb{Z} \to
\operatorname{GL}_2(\mathbb{C})$
representasi dengan
\[
\sigma(1) = \begin{pmatrix} 0&-1\\1&0 \end{pmatrix} .
\]
Definisikan $\rho'(g)=\rho(g^{-1})$
dan $\sigma'(g)=\sigma(g^{-1})$ untuk semua $g \in G$.
Dengan mempertimbangkan nilai eigen, kita menyimpulkan bahwa
$\sigma'$ dan $\sigma$ ekuivalen, sedangkan $\rho'$ dan $\rho$ tidak.
\end{example}

\begin{example} \label{ex:S3_decomposition}
Misalkan $\rho$ representasi permutasi berdimensi $3$ yang terkait dengan
aksi $S_3$ pada $\{1,2,3\}$.
Dalam basis biasa, kita mempunyai
\[
\rho(23) = \begin{pmatrix} 1&0&0\\ 0&0&1 \\ 0&1&0 \end{pmatrix}
\textrm{ dan }
\rho(123) = \begin{pmatrix} 0&0&1\\1&0&0\\0&1&0 \end{pmatrix}.
\]
Andaikan $\zeta \in k$ suatu akar satuan kubik primitif.
Melalui matriks perubahan basis
\[
P = \begin{pmatrix} 1&1&1\\ 1&\zeta^2&\zeta\\1&\zeta&\zeta^2 \end{pmatrix} 
\]
kita memperoleh representasi baru $\sigma$ dengan
\[
\sigma(23) = \begin{pmatrix} 1&0&0\\ 0&0&1 \\ 0&1&0 \end{pmatrix}
\textrm{ dan }
\sigma(123) = \begin{pmatrix} 1&0&0\\ 0&\zeta&0 \\ 0&0&\zeta^2
\end{pmatrix} .  
\]
Perhatikan bahwa $\rho$ dan $\sigma$ serupa berdasarkan konstruksi.
\end{example}


\subsection{Subrepresentasi}

\begin{definition}
Misalkan $\rho : G \to \operatorname{GL}(V)$ suatu representasi linear.
Suatu subruang $W \subset V$ disebut \emph{stabil di bawah aksi $G$} jika
$\rho_g(w) \in W$ untuk semua $g \in G$, $w \in W$.
Definisikan $\rho^W : G \to \operatorname{GL}(W)$ dengan $\rho^W(g) := \rho(g)|_W$ untuk semua $g
\in G$.
Kita menyebut $\rho^W$ suatu \emph{subrepresentasi} dari $\rho$.
\end{definition}

Dalam situasi di atas, kita juga cukup mengatakan “$W$ adalah subrepresentasi dari $V$”
saja.

Jika $W$ merupakan subrepresentasi dari $V$, maka
\[
\rho_g(v+w) = \rho_g(v) + \rho_g(w) \in \rho_g(v)+W
\]
untuk setiap $v \in V$, $w \in W$, dan $g \in G$.
Hal ini menghasilkan \emph{representasi hasil bagi} yang terdefinisi dengan baik,
$\rho' : G \to \operatorname{GL}(V/W)$,
pada ruang hasil bagi $V/W$.

\begin{definition}
Untuk representasi $(V,\rho)$ dan $(W,\sigma)$ dari suatu grup $G$,
\emph{representasi jumlah langsung} $\rho \oplus \sigma$
adalah representasi dengan ruang dasar $V \oplus W$ yang diberikan oleh
\[ (\rho \oplus \sigma)_g (v,w) := ( \rho_g(v), \sigma_g(v) ) \]
untuk semua $v \in V$, $w \in W$, dan $g \in G$.
\end{definition}

\begin{proposition}
Misalkan $V$ dan $W$ representasi dari $G$, dan misalkan
$f : V \to W$ suatu pemetaan $G$-ekuivarian.
Maka $\ker(f)$ merupakan subrepresentasi dari $V$.
\end{proposition}

\begin{example}
Dalam notasi
Contoh~\ref{ex:S3_decomposition}, kita melihat bahwa $V$ dapat ditulis
sebagai jumlah langsung $U \oplus W$, dengan
\[
U = \operatorname{span}\left\{ \begin{pmatrix} 1,1,1 \end{pmatrix}^T
\right\}
\]
merupakan representasi trivial, dan
\[
W = \operatorname{span}\left\{
\begin{pmatrix}1,\zeta,\zeta^2\end{pmatrix}^T,
\begin{pmatrix}1,\zeta^2,\zeta\end{pmatrix}^T
\right\}
\]
merupakan representasi setia berdimensi $2$ dari $S_3$.
\end{example}

\begin{example}
Andaikan $G=\langle g \rangle$ suatu grup siklik hingga.
Ingat bahwa setiap matriks berorde hingga dapat didiagonalkan
di atas $\mathbb{C}$.
Jadi, setiap representasi kompleks berdimensi hingga dari
$G$ merupakan jumlah langsung representasi-representasi berdimensi $1$.
\end{example}

\begin{example}
Jika $V$ suatu representasi dari $G$, definisikan \emph{ruang invarian}
\[
V^G = \{ v \in V \mid gv=v \textrm{ untuk semua } g \in G \}.
\]
Perhatikan bahwa $V^G$ merupakan subrepresentasi dari $V$ dan merupakan jumlah langsung
representasi-representasi trivial.
\end{example}

\begin{definition}
Representasi linear $V$ disebut \emph{tak tereduksi} jika $V \ne 0$ dan satu-satunya
subrepresentasi dari $V$ adalah $0$ dan $V$ sendiri.
Jika tidak demikian, representasi tersebut disebut \emph{tereduksi}.
\end{definition}

\begin{definition}
Representasi linear $V$ disebut \emph{tak terdekomposisi} jika $V \ne 0$ dan
$V$ tidak dapat ditulis sebagai jumlah langsung $V = W_1 \oplus W_2$ dengan
$W_1, W_2$ subrepresentasi tak nol.
Jika tidak demikian, representasi tersebut disebut \emph{terdekomposisi}.
\end{definition}

Representasi terdekomposisi jelas tereduksi, tetapi kebalikannya
belum tentu berlaku, seperti diperlihatkan oleh contoh berikut.

\begin{example}
Misalkan $G=\mathbb{Z}/p\mathbb{Z}$ untuk suatu bilangan prima $p$,
dan pertimbangkan representasi berdimensi dua $\rho$ di atas $\mathbb{F}_p$
yang diberikan oleh
\[
\rho(1) = \begin{pmatrix} 1&1\\0&1 \end{pmatrix} .
\]
Suatu vektor $\begin{pmatrix} a\\b \end{pmatrix} $ merentang
subruang stabil-$G$ hanya jika $b=0$.
Subruang yang direntang oleh $\begin{pmatrix} 1\\0 \end{pmatrix}$
merupakan subrepresentasi tak nol sejati
dari $\rho$, sehingga $\rho$ tereduksi.
Namun, inilah satu-satunya subrepresentasi, sehingga $\rho$
tidak dapat ditulis sebagai jumlah langsung subrepresentasi-subrepresentasi sejati.
Kita menyimpulkan bahwa $\rho$ tak terdekomposisi.
\end{example}

\begin{theorem}[Krull--Schmidt]
Andaikan $V$ suatu representasi berdimensi hingga dari grup hingga $G$.
Semua dekomposisi
\[
V \cong W_1 \oplus W_2 \oplus \cdots \oplus W_r
\]
menjadi representasi-representasi tak terdekomposisi bersifat tunggal hingga isomorfisme
dan pengurutan ulang.
\end{theorem}

\begin{proof}
Untuk sementara ditunda. Kita akan membuktikan hasil ini dalam keumuman yang lebih besar ketika membahas
modul.
\end{proof}

\begin{remark}
Teorema Krull--Schmidt gagal untuk representasi $G \to
\operatorname{GL}(\mathbb{Z})$, sehingga fakta bahwa kita bekerja
di atas medan sangat penting di sini.
Sayangnya, contoh-contoh tandingnya cukup halus!
\end{remark}

\begin{theorem}[Lema Schur]
Jika $V$ dan $W$ merupakan representasi tak tereduksi berdimensi hingga
dan $k$ tertutup secara aljabar, maka
\[
\operatorname{Hom}_k^G(V,W) =
\begin{cases}
0 & \textrm{ jika } V \not\cong W\\
k & \textrm{ jika } V \cong W.
\end{cases}
\]
\end{theorem}

\begin{proof}
Misalkan $\phi : V \to W$ suatu pemetaan $G$-ekuivarian.
Karena $V$ tak tereduksi, berlaku $\ker(\phi)=0$ atau $\ker(\phi)=V$.
Karena $W$ tak tereduksi, berlaku $\operatorname{im}(\phi)=0$
atau $\operatorname{im}(\phi)=W$.
Jadi, $\phi$ adalah $0$ atau suatu isomorfisme.

Dengan demikian, kita direduksi ke kasus $W=V$.
Sekarang $\operatorname{End}^G_k(V)=\operatorname{Hom}_k^G(V,W)$
merupakan aljabar-$k$ terhadap komposisi.
Jika $\phi \in \operatorname{End}^G_k(V)$ tak nol, maka ia harus merupakan isomorfisme.
Jadi, $\phi^{-1}$ juga termasuk dalam $\operatorname{End}^G_k(V)$.
Kita menyimpulkan bahwa $\operatorname{End}^G_k(V)$ merupakan aljabar pembagian.
Hasil tersebut sekarang mengikuti dari Lema~\ref{lem:division_over_kbar}.
\end{proof}

\begin{lemma} \label{lem:division_over_kbar}
Jika $k$ suatu medan tertutup secara aljabar,
maka setiap aljabar pembagian-$k$ berdimensi hingga $D$
isomorfik dengan $k$. 
\end{lemma}

\begin{proof}
Andaikan $a \in D \setminus \{0\} $.
Karena $D$ berdimensi hingga, perkalian dengan $a$
mempunyai polinomial minimal $m_a(t) \in k[t]$.
Jika $m_a$ tidak berderajat $1$, maka $m_a(t)=f(t)g(t)$
untuk polinomial tak konstan $f,g \in k[t]$, sebab $k$ tertutup secara
aljabar.
Persamaan $0=m_a(a)=f(a)g(a)$ memaksa $f(a)=0$ atau $g(a)=0$,
sebab $D$ merupakan gelanggang pembagian. Namun, ini bertentangan dengan keminimalan
polinomial minimal. Kita menyimpulkan bahwa $m_a(t)=t-\lambda$ untuk suatu
$\lambda \in k$. Jadi, $a$ merupakan perkalian skalar dengan $\lambda$,
dan dengan demikian merupakan elemen $k$.
\end{proof}

\begin{definition}
Representasi linear $V$ disebut \emph{tereduksi lengkap}
jika $V$ merupakan jumlah langsung representasi-representasi tak tereduksi.
\end{definition}

Andaikan $V$ tereduksi lengkap.
Teorema Krull--Schmidt menyatakan bahwa terdapat
representasi-representasi tak tereduksi $W_1, \ldots, W_r$ yang saling tak isomorfik,
sedemikian sehingga
\begin{equation} \label{eq:isotypic_decomposition}
V = \bigoplus_{i=1}^r V_i
\end{equation}
dengan setiap $V_i$ isomorfik dengan
\[
V_i = W_i^{\oplus m_i}
\]
untuk suatu bilangan bulat positif $m_i$.
Pernyataan \eqref{eq:isotypic_decomposition} disebut
\emph{dekomposisi isotipik dari $V$}.
Subrepresentasi $V_i$ disebut \emph{komponen isotipik dari
$V$ yang terkait dengan $W_i$}, sedangkan bilangan bulat $m_i$ disebut
\emph{multiplisitas $W_i$ dalam $V$}.

Setiap $V_i$ merupakan subrepresentasi \emph{kanonik} dari $V$,
sedangkan klasifikasi subrepresentasi yang isomorfik dengan $W_i$ dapat bergantung pada
pilihan sebarang.
Hal ini memperumum keadaan ketika ruang-ruang eigen dari suatu transformasi linear
bersifat kanonik, sedangkan basis yang terdiri atas vektor-vektor eigen
dapat bergantung pada pilihan sebarang.

\begin{exercise}
Buktikan Teorema Krull--Schmidt untuk kasus representasi
tereduksi lengkap $V$ ketika medan $k$ tertutup secara aljabar.
(Petunjuk: gunakan Lema Schur.)
\end{exercise}

Teorema berikut merupakan alasan utama bahwa teori representasi grup hingga
dalam karakteristik “baik” sangat berbeda dari teori dalam
karakteristik “buruk”.
Berdasarkan konvensi, setiap bilangan bulat dianggap “relatif prima terhadap karakteristik $k$”
ketika karakteristiknya $0$.

\begin{theorem}[Maschke]
Misalkan $G$ suatu grup hingga yang ordenya relatif prima terhadap karakteristik $k$.
Setiap representasi berdimensi hingga dari $G$ tereduksi lengkap.
\end{theorem}

Teorema ini mengikuti melalui induksi dan lema berikut:

\begin{lemma}
Misalkan $G$ suatu grup hingga yang ordenya relatif prima terhadap karakteristik $k$.
Jika $V$ suatu representasi dari $G$ dan $W$ subrepresentasi dari $V$,
maka terdapat subrepresentasi $U$ dari $V$
sedemikian sehingga $V= W \oplus U$.
\end{lemma}

\begin{proof}
Misalkan $Q: V \to W$ suatu proyeksi dari $V$ ke subruang $W$.
Proyeksi $Q$ ada berdasarkan alasan aljabar linear
dan belum tentu \emph{$G$-ekuivarian}.
Namun, kita akan menyesuaikannya agar \emph{menjadi} $G$-ekuivarian.
Definisikan $P : V \to W$ melalui
\[
P(v) = \frac{1}{|G|} \sum_{g \in G} g\left( Q\left( g^{-1} v\right)\right)
\]
untuk semua $v \in V$. Perhatikan bahwa syarat pada karakteristik
medan diperlukan agar kita dapat membagi dengan orde grup.

Kita mengklaim bahwa $P : V \to W$ merupakan proyeksi $G$-ekuivarian ke $W$.
Karena pembatasan $Q$ pada $W$ adalah identitas dan $W$ stabil-$G$,
kita melihat bahwa
$g \circ Q \circ g^{-1}$ merupakan proyeksi ke $W$
untuk semua $g \in G$. Jadi, kita mempunyai
\begin{align*}
P^2(v) &=
\frac{1}{|G|^2} \sum_{g,h \in G} gQg^{-1} hQh^{-1}v\\
&=\frac{1}{|G|^2} \sum_{g,h \in G} hQh^{-1}v\\
&=\frac{1}{|G|} \sum_{h \in G} hQh^{-1}v = P(v)
\end{align*}
dan menyimpulkan bahwa $P$ merupakan proyeksi. Karena
$P|_W=\operatorname{id}_W$, kita menyimpulkan bahwa ia merupakan proyeksi ke $W$.

Sekarang kita tunjukkan bahwa $P$ ekuivarian.
Hal ini mengikuti dari kiat pengindeksan ulang $j=h^{-1}g$ berikut:
\begin{align*}
P(hv) &=
\frac{1}{|G|} \sum_{g \in G} gQg^{-1} (hv)\\
&=\frac{1}{|G|^2} \sum_{j \in G} hjQj^{-1}v\\
&=h\left(\frac{1}{|G|} \sum_{j \in G} jQj^{-1}v\right) = hP(v).
\end{align*}
Kernel $P$ merupakan komplemen stabil-$G$, $U$, yang dicari.
\end{proof}

\subsection{Konstruksi}

Andaikan $(V,\rho)$ dan $(W,\sigma)$ representasi linear berdimensi hingga
dari suatu grup $G$ di atas medan $k$.

Misalkan $\operatorname{Hom}_k(V,W)$ menyatakan ruang vektor transformasi linear
$\tau : V \to W$.
Misalkan $V^\vee$ ruang vektor dual dari $V$.
Misalkan $V \otimes W$ hasil kali tensor di atas $k$.

\begin{proposition}
Ruang vektor $\operatorname{Hom}_k(V,W)$ mempunyai struktur kanonik sebagai representasi
linear $\tau$, dengan
\[
 \tau_g(f) = \sigma(g) \circ f \circ \rho(g)^{-1}
\]
untuk semua $f : V \to W$ dan $g \in G$.
\end{proposition}

Aksi-$G$ yang didefinisikan di atas mempunyai sifat sangat berguna bahwa
himpunan homomorfisme $G$-ekuivarian tepat merupakan himpunan invarian dari
aksi-$G$ pada himpunan homomorfisme. Dengan kata lain,
\[
\operatorname{Hom}^G_k(V,W) = \operatorname{Hom}_k(V,W)^G .
\]

\begin{proposition}
Ruang dual $V^\vee$ mempunyai struktur kanonik sebagai representasi
linear $\rho^\vee$, dengan
\[
 (\rho^\vee)_g(f) = f \circ \rho(g)^{-1}
\]
untuk $f : V \to k$ dan $g \in G$.
\end{proposition}

Perhatikan bahwa proposisi di atas merupakan kasus khusus aksi pada
homomorfisme ketika $W=k$ mempunyai aksi-$G$ trivial.

\begin{proposition}
Hasil kali tensor $V \otimes W$ mempunyai struktur kanonik sebagai representasi
linear $\rho \otimes \sigma$, dengan
\[
(\rho \otimes \sigma)_g(v \otimes w) = \rho_g(v) \otimes \sigma_g(w)
\]
untuk $v \in V$, $w \in W$, dan $g \in G$.
\end{proposition}

Perhatikan bahwa proposisi di atas sesuai dengan aksi-$G$ yang
diperoleh dari isomorfisme kanonik
\[
\operatorname{Hom}_k(W,V) \cong W^\vee \otimes V. 
\]

\begin{exercise}
Untuk suatu ruang vektor $V$,
tentukan aksi $G$ pada
$\mathcal{T}^d(V)$, $\mathcal{S}^d(V)$, $\Lambda^d(V)$,
$\operatorname{Sym}^n(V)$, dan $\operatorname{Alt}^n(V)$.
Periksa bahwa berbagai pemetaan kanonik bersifat ekuivarian.
\end{exercise}

%%%%%%%%%%%%%%%%%%
%%%%%%%%%%%%%%%%%%

\bibliographystyle{alpha}
\bibliography{rep_theory}

\end{document}


